\textbf{1.} Pour $n = 1$, l'inégalité s'écrit $\sqrt{ab} \leqslant \frac{a+b}{2}$ :
c'est l'inégalité entre moyennes géométrique et arithmétique de deux nombres.

\textbf{2.} Les $n+1$ nombres $a, \underbrace{b, b, \ldots, b}_{n \text{ fois}}$
conviennent : leur produit est $ab^n$ et leur somme est $a + nb$.

\textbf{3.} L'inégalité arithmético-géométrique pour ces $n+1$ nombres
strictement positifs donne
\[
\sqrt[n+1]{a \cdot b \cdots b} \leqslant \frac{a + b + \cdots + b}{n+1},
\qquad\text{soit}\qquad
\sqrt[n+1]{ab^n} \leqslant \frac{a + nb}{n+1}.
\]
L'égalité a lieu si et seulement si les $n+1$ nombres sont égaux,
c'est-à-dire si et seulement si $a = b$.

\textbf{4.} Avec $a = 1$ et $b = 1 + \frac1n$, on a $a \neq b$, donc
l'inégalité est stricte :
\[
\sqrt[n+1]{\Bigl(1 + \frac1n\Bigr)^n} < \frac{1 + n + 1}{n+1} = 1 + \frac{1}{n+1}.
\]
Les deux membres sont positifs ; en les élevant à la puissance $n+1$
(fonction strictement croissante sur $[0, +\infty[$), on obtient
$\left(1 + \frac1n\right)^n < \left(1 + \frac1{n+1}\right)^{n+1}$.
La suite de terme général $\left(1 + \frac1n\right)^n$ est donc strictement
croissante. \hfill $\blacksquare$
